\documentclass[10pt]{article}
\usepackage[a4paper,margin=1in]{geometry}
\usepackage{xeCJK}
\setCJKmainfont{FandolSong-Regular.otf}[ItalicFont=FandolKai-Regular.otf]
\setCJKsansfont{FandolHei-Regular.otf}
\setCJKmonofont{FandolFang-Regular.otf}
\usepackage{hyperref}
\usepackage{enumitem}
\setlist{nosep,leftmargin=1.5em}
\hypersetup{colorlinks=true,linkcolor=blue,urlcolor=blue,citecolor=blue}
\XeTeXlinebreaklocale "zh"
\XeTeXlinebreakskip = 0pt plus 1pt
\setlength{\parindent}{2em}
\setlength{\parskip}{0.35em}
\emergencystretch=3em
\sloppy

\title{Temporal Difference Learning for Model Predictive Control\\[0.5em]\large 中文重构译文}
\author{中文译文备份}
\date{}

\begin{document}
\maketitle

\noindent\textbf{说明：} 本文档由 PDF 抽取文本重构生成，用于在缺少可用 arXiv LaTeX 源码时备份中文译文。章节和段落为自动恢复，图表、公式和双栏排版未按原始论文完整复刻。

\bigskip

\section*{Temporal Difference Learning for Model Predictive Control}

Nicklas Hansen，Xiaolong Wang，Hao Su

arXiv:2203.04955v2 [cs.LG]，2022年7月19日。

\section*{摘要}
数据驱动的模型预测控制相对于无模型方法具有两个关键优势：一是通过学习模型提高样本效率，二是随着规划计算预算增加，性能也可以继续提升。本文结合无模型方法和基于模型方法的优势，使用面向任务的潜在动力学模型进行短时域局部轨迹优化，并使用学得的终端价值函数估计长期回报；二者都通过时间差分学习共同训练。

本文方法 TD-MPC 在状态和图像输入的连续控制任务上取得了优于既有模型方法和无模型方法的样本效率与最终性能。代码和视频见：\url{https://nicklashansen.github.io/td-mpc}。

图1. 概览。上方：TD-MPC 使用面向任务的潜在动力学模型和学得的价值函数进行模型预测控制。方法在模型 rollout 上优化轨迹，并利用价值函数增强长期回报估计。下方：在人形和四足机器人运动任务上，TD-MPC 相比 SAC 以及使用真实模拟器的 MPC 取得更高 episode return；阴影区域表示 95\% 置信区间。

\section{导言}
强化学习智能体通常需要通过反复交互来获取环境知识。规划是一类重要的序贯决策方法：智能体先借助内部模型评估候选动作序列，再执行其中的第一个动作。与完全依靠试错来学习策略的无模型算法相比，基于模型的方法有机会利用模型预测提高样本效率；但在连续控制中，模型误差容易沿着长时域 rollout 累积，长时域规划本身也很昂贵。

已有基于模型的工作大致可以分成两条路线。第一类直接用模型做规划，利用轨迹优化得到动作；当规划时域变长时，这类方法的计算代价迅速上升。第二类把学得模型生成的 rollout 用于改进无模型策略，样本效率通常更好，但模型偏差会传递到策略和价值估计中。这也是基于模型方法在若干连续控制任务上长期难以稳定超过强无模型基线的原因。

本文提出的 TD-MPC（Temporal Difference Learning for Model Predictive Control）试图结合两者的长处。它使用一个面向任务的潜在动力学模型进行短时域模型预测控制，同时使用通过 TD 学习得到的终端价值函数估计规划时域之外的长期回报。每个决策步中，TD-MPC 根据学得模型给出的短期奖励和价值函数给出的长期估计来优化动作序列，只执行序列中的第一个动作，然后在下一步重新规划。

TD-MPC 的关键在于模型学习方式。作者没有要求模型重建完整状态或图像，因为阴影、背景纹理等细节通常与控制目标无关，却会增加预测难度。相反，模型只学习与奖励和价值估计有关的潜在表示，并把奖励预测损失、TD 价值损失以及潜在状态一致性损失一起反向传播到多步模型 rollout 中。这样既保留了规划所需的局部动力学，又减少了对观测模态和像素级细节的依赖。

作者在 DMControl 和 Meta-World 的连续控制任务上评估该方法。结果显示，TD-MPC 在状态输入和图像输入任务中都能达到较好的样本效率，并能处理 Humanoid 与 Dog 等高维动作空间任务。对于图 1 中的人形运动任务，可以把底层关节协调交给短时域模型规划，把跑动方向等长期目标交给价值函数估计，两者共同决定最终动作。

\section{预备知识}
\subsection*{问题设定}
论文考虑无限时域马尔可夫决策过程，记为 $(S,A,T,R,\gamma,p_0)$。其中 $S$ 和 $A$ 分别表示连续状态空间和连续动作空间，$T$ 是动力学转移函数，$R$ 是奖励函数，$\gamma$ 是折扣因子，$p_0$ 是初始状态分布。目标是学习一个策略或动作选择映射，使沿轨迹获得的折扣回报最大。

\subsection*{拟合 Q 迭代}
无模型 TD 学习通常估计最优状态-动作价值函数 $Q^\ast(s,a)$。在参数化情形下，算法维护 $Q_\theta(s,a)$，并用重放缓冲区中的转移样本拟合 TD 目标：
\[
\theta_{k+1}\leftarrow \arg\min_\theta E_{(s,a,s')\sim B}\|Q_\theta(s,a)-y\|_2^2,\quad
y=R(s,a)+\gamma\max_{a'}Q_{\theta^-}(s',a').
\]
这里 $B$ 是重放缓冲区，$\theta^-$ 是由在线参数 $\theta$ 的慢速移动平均得到的目标网络参数。若 $\gamma$ 接近 1，$Q_\theta$ 近似描述从当前状态-动作对出发的长期最优回报。

\subsection*{模型预测控制}
在传统控制中，动作选择常被写成一个有限时域轨迹优化问题。MPC 在当前状态 $s_t$ 下优化未来 $H$ 步动作序列：
\[
\Pi_\theta^{\mathrm{MPC}}(s_t)=
\arg\max_{a_{t:t+H}} E\left[\sum_{i=t}^{t+H}\gamma^i R(s_i,a_i)\right].
\]
实际执行时只采用第一个动作 $a_t$，下一步再重新优化。这种滚动时域方式计算可控，但只看有限时域，天然偏近视。若能得到终端价值函数，就可以把 $s_{t+H}$ 之后的长期回报纳入估计，这正是 TD-MPC 使用 TD 学习补足 MPC 的位置。

\section{用于模型预测控制的 TD 学习}
TD-MPC 在推理阶段使用 MPPI（Model Predictive Path Integral）做采样式轨迹优化。给定当前观测 $s_t$，表示网络先将其编码为潜在状态 $z_t=h_\theta(s_t)$。随后，规划器在潜在空间中用动力学模型 $d_\theta$ 推进状态，用奖励模型 $R_\theta$ 估计短期奖励，并用终端价值函数 $Q_\theta$ 估计 rollout 末端之后的长期回报。对一条候选轨迹 $\Gamma$，其估计回报由模型 rollout 中的奖励项和末端价值项共同组成。

MPPI 维护一个关于动作序列的高斯采样分布。每轮迭代中，规划器采样多条长度为 $H$ 的候选轨迹，计算它们的估计回报，选取回报最高的一批轨迹，并按回报重新加权更新均值与方差。固定迭代次数后，从最终分布中采样动作序列。由于 MPC 采用滚动时域控制，TD-MPC 在每个决策步只执行第一个动作，并把上一步得到的均值右移一位作为下一次规划的 warm start，从而减少收敛所需迭代次数。

探索也通过规划过程完成。单纯依赖采样方差会导致不同任务上的探索程度差异很大，因此 TD-MPC 对动作分布的标准差设置随训练线性衰减的下界。训练初期模型尚不准确，方法还会把规划时域从 1 逐步增加到最终的 $H$，避免早期规划过多依赖不可靠的多步模型预测。

除采样轨迹外，TD-MPC 还学习一个确定性策略 $\pi_\theta$ 来辅助规划。规划器把由 $\pi_\theta$ 产生的轨迹加入候选集合：如果策略轨迹估计回报低，它会被 top-k 筛选自然排除；如果估计回报高，它会按回报权重影响分布更新。训练时，策略动作上加入线性衰减的高斯噪声，使策略引导不至于过早收缩探索。

\section{面向任务的潜在动力学模型}
TD-MPC 使用的 TOLD（Task-Oriented Latent Dynamics）模型只刻画与任务有关的动力学，而不是复原整个环境。模型由五个学习组件组成：
\begin{itemize}
  \item 表示网络：$z_t=h_\theta(s_t)$；
  \item 潜在动力学：$z_{t+1}=d_\theta(z_t,a_t)$；
  \item 奖励预测：$\hat r_t=R_\theta(z_t,a_t)$；
  \item 价值预测：$\hat q_t=Q_\theta(z_t,a_t)$；
  \item 策略网络：$\hat a_t=\pi_\theta(z_t)$。
\end{itemize}

给定时刻 $t$ 的观测，TOLD 只编码第一帧或第一个状态，然后在潜在空间中递归预测后续潜在状态。每一步都会预测奖励、价值和策略动作。所有组件均可用确定性 MLP 实现，不需要循环门控结构或概率状态空间模型。这个设计牺牲了对环境细节的完整复原能力，但更贴合控制任务需要。

训练目标是一个时间加权的多步损失。单步损失包含三项：
\[
L_i=
c_1\|R_\theta(z_i,a_i)-r_i\|_2^2
+c_2\|Q_\theta(z_i,a_i)-(r_i+\gamma Q_{\theta^-}(z_{i+1},\pi_\theta(z_{i+1})))\|_2^2
+c_3\|d_\theta(z_i,a_i)-h_{\theta^-}(s_{i+1})\|_2^2.
\]
第一项拟合即时奖励，第二项是 TD 价值目标，第三项要求模型预测出的下一潜在状态接近目标网络编码出的真实下一观测表示。目标网络参数 $\theta^-$ 由在线参数的指数移动平均得到。

潜在状态一致性是 TOLD 的核心约束。早期世界模型通常通过预测未来状态或像素来学习动力学，但这会迫使模型解释大量与任务无关的观测细节。TOLD 只要求预测的潜在状态与真实下一观测的潜在表示一致，因此可以适配状态、图像和多模态输入。由于梯度会从奖励、价值和一致性三类损失穿过多个 rollout 步反传，模型学习到的潜在空间更直接服务于规划和价值估计。

训练过程在数据收集和模型更新之间交替进行。智能体用当前 TD-MPC 规划器与环境交互，把转移样本写入重放缓冲区；随后从缓冲区采样长度为 $H$ 的轨迹，编码首个观测，在潜在空间展开 TOLD，并用上述损失更新表示、动力学、奖励、价值和策略网络。策略 $\pi_\theta$ 还用于近似 $Q$ 函数的最大化动作，避免在每次计算 TD 目标时都运行昂贵的规划过程。

图 2 可理解为这一训练流程的示意：从重放缓冲区采样一段轨迹，首个观测被编码为 $z_0$，TOLD 递归预测 $z_1,\ldots,z_H$，同时预测每个潜在状态上的奖励、价值和动作；后续真实观测仅通过目标网络编码后作为潜在一致性目标，不要求模型重建原始观测。

\section{实验设置与主要结果}
作者在 92 个连续控制任务上评估 TD-MPC，任务来自 DeepMind Control Suite 和 Meta-World v2，覆盖稀疏奖励、高维状态与动作空间、图像观测、多模态输入、目标条件控制和多任务学习。论文主要回答三个问题：TD-MPC 与强无模型和基于模型方法相比表现如何；TOLD 在以奖励为中心的目标下是否仍具备多任务和迁移能力；规划预算与性能之间有什么关系。

实现上，所有 TOLD 组件均为确定性网络。状态输入任务主要使用 MLP，图像输入任务使用卷积编码器。默认规划时域为 $H=5$；规划时通常采样 $N=512$ 条轨迹，并把约 5\% 的候选轨迹来自策略网络。规划迭代次数默认较少，Dog 和 Humanoid 等更复杂任务使用更多迭代。作者报告该实现可在单张 GPU 上较快求解大多数任务。

比较方法包括 SAC、LOOP、使用真实模拟器规划的 MPC、CURL、DrQ、DrQ-v2、PlaNet、Dreamer、DreamerV2、MuZero 和 EfficientZero。消融实验比较了无潜在表示、无潜在一致性损失，以及用重建目标或对比目标替代一致性损失的版本。

在基于状态的 DMControl 任务中，TD-MPC 在 Quadruped、Acrobot、Humanoid 和 Dog 等动力学复杂或动作维度较高的任务上优势更明显。使用真实模拟器的 MPC 在部分运动任务上表现不错，但在稀疏奖励任务中容易失败。无模型 SAC 是强基线，但在高维复杂任务上样本效率通常不及 TD-MPC。

在基于图像的 DMControl 100k 基准中，TD-MPC 没有专门为图像强化学习调参，却能与 DrQ、DrQ-v2、DreamerV2 等方法竞争。MuZero 和 EfficientZero 依赖离散化动作空间，在高维连续动作任务上会遇到维度爆炸，而 TD-MPC 的规划过程天然支持连续动作。

\begin{center}
\scriptsize
\begin{tabular}{lrrrrrrrrr}
任务 & SAC 状态 & SAC 像素 & CURL & DrQ & PlaNet & Dreamer & MuZero & Eff.Zero & TD-MPC \\
Cartpole Swingup & $812\pm45$ & $419\pm40$ & $597\pm170$ & $759\pm92$ & $563\pm73$ & $326\pm27$ & $219\pm122$ & $813\pm19$ & $770\pm70$ \\
Reacher Easy & $919\pm123$ & $145\pm30$ & $517\pm113$ & $601\pm213$ & $82\pm174$ & $314\pm155$ & $493\pm145$ & $952\pm34$ & $628\pm105$ \\
Cup Catch & $957\pm26$ & $312\pm63$ & $772\pm241$ & $913\pm53$ & $710\pm217$ & $246\pm174$ & $542\pm270$ & $942\pm17$ & $933\pm24$ \\
Finger Spin & $672\pm76$ & $166\pm128$ & $779\pm108$ & $901\pm104$ & $560\pm77$ & $341\pm70$ & -- & -- & $943\pm59$ \\
Walker Walk & $604\pm317$ & $42\pm12$ & $344\pm132$ & $612\pm164$ & $221\pm43$ & $277\pm12$ & -- & -- & $577\pm208$ \\
Cheetah Run & $228\pm95$ & $103\pm38$ & $307\pm48$ & $344\pm67$ & $165\pm123$ & $235\pm137$ & -- & -- & $222\pm88$ \\
\end{tabular}
\end{center}

表 1 汇总了图像输入 100k 环境步基准上的回报。星号方法在原文中使用离散动作空间；这里保留核心数值，去除了 PDF 抽取时混入的无关文本。

在 Meta-World、目标条件控制和多任务实验中，TD-MPC 也能保持较好的样本效率。多模态实验表明，把本体状态与第一视角相机输入一起使用时，TOLD 的潜在表示能够融合不同输入源。可变计算预算实验则显示，增加规划时域或 MPPI 迭代次数通常能提升性能，但默认预算已经在计算成本与回报之间取得了较好的折中；推理时也可以只使用共同学习到的策略，从而降低规划开销。

训练耗时方面，论文在单张 RTX 3090 GPU 上比较 SAC、LOOP、MPC:sim 和 TD-MPC。TD-MPC 在 Walker Walk 上比 LOOP 更快达到目标回报，同时每 50 万环境步的训练成本更低。作者据此认为，TD-MPC 在样本效率和实际计算成本之间比若干混合式基线更均衡。

\section{相关工作}
TD-MPC 与无模型 TD 学习、基于模型强化学习和混合式方法都有联系。DDPG、SAC 等 actor-critic 方法学习策略和价值函数，但不学习环境模型；PlaNet、Dreamer 等方法学习潜在动力学，通常依赖重建或图像预测目标；MuZero 和 EfficientZero 则从奖励和价值信号学习潜在模型，并通过 MCTS 做离散动作选择。

TD-MPC 的区别在于，它同时学习面向任务的潜在模型、价值函数和策略，并在推理时用采样式 MPC 选择连续动作。与需要像素重建的世界模型相比，TOLD 更强调奖励、价值和潜在一致性；与纯无模型方法相比，TD-MPC 在动作选择阶段仍保留了规划带来的计算可扩展性。

\section{结论}
论文表明，把 TD 学习得到的终端价值函数接入短时域 MPC，可以缓解有限时域规划的近视问题；而面向任务的潜在动力学模型又避免了完整重建环境带来的负担。TD-MPC 因此在一系列连续控制任务上兼具样本效率和较强的最终性能，并能自然扩展到图像输入、多模态输入和高维连续动作空间。

作者也指出，后续工作可以继续改进探索机制、模型使用方式和规划效率。对个人复现而言，主要门槛不在数据集规模，而在于如何用较少环境步训练紧凑潜在模型，并为每个决策步预留一定规划计算。


\section*{参考文献与附录说明}
原始 PDF 抽取出的参考文献和附录包含大量双栏错位、公式残片和乱码，逐条保留会明显影响可读性。本重构版暂不逐条复刻参考文献排版，而是保留正文中的作者年份引用，并把附录内容整理为要点，便于后续继续精修。

\subsection*{附录要点概述}
附录主要补充四类内容。第一，作者考察了 TOLD 模型在相关任务之间的迁移能力：先在 Walk 类任务上训练，再迁移到 Run 类任务；结果显示，表示网络 $h_\theta$ 具有一定可迁移性，而动力学预测器 $d_\theta$ 更偏向任务特定行为。第二，论文把 TD-MPC 与 SAC、LOOP、PlaNet、Dreamer、MuZero、EfficientZero 等方法在模型目标、推理方式、连续动作支持和计算成本上做了定性比较。第三，附录讨论了潜在一致性损失的替代方案，包括对比损失和重建损失；实验支持作者采用的潜在一致性目标。第四，附录列出了实现细节、超参数、网络结构、基线设置、可变规划预算实验、聚合指标和更多 Meta-World 任务可视化。

从复现角度看，附录中最值得保留的是默认超参数：规划时域通常为 $H=5$，MPPI 迭代次数约为 6，复杂 Dog 与 Humanoid 任务使用更多迭代；采样轨迹数约为 512；状态任务主要使用 MLP，图像任务使用卷积编码器；重放缓冲区采用优先经验回放，优先级与价值损失相关。论文还强调，训练和评估中的“环境步”与“决策步”需要区分，前者不依赖动作重复，后者与策略查询次数相关。

\end{document}
